% !TEX TS-program = xelatex
\documentclass[12pt,letterpaper]{article}
\usepackage{style/dsc180reportstyle}
\title{BaZiReason: Evaluating Tool-Augmented LLM Reliability in Traditional Chinese Calendar Reasoning}
\author{Yonghao Wang\\DSC 180A, Fall 2026\\Domain D25 --- Advisor: Jingbo Shang}
\date{}
\setlength{\emergencystretch}{2em}
\begin{document}
\maketitle
\begin{abstract}
Chinese Four Pillars (BaZi) chart construction combines sexagenary cycles, solar-term transitions, and time conventions. This provides a constrained setting for studying language-model reliability without evaluating claims about destiny. BaZiReason asks whether deterministic context improves calendar-answer reliability, and which gains reflect preprocessing, faithful copying, or responses to explicit uncertainty policy. An independent audit of the initial 96-case checkpoint reproduced its software tests but identified duplicate inputs, incomplete boundary coverage, timezone-cache dependence, and insufficient experimental controls. The revised suite contains 152 synthetic cases, including all twelve month-start solar terms and matched missing-information cases. A separate Skyfield/JPL computation checks 24 solar-term timestamps: lunar-python differs by $-1.17$ to $+2.26$ seconds, whereas sxtwl is approximately 13--16 seconds earlier. These checks strengthen timestamp evidence without establishing complete-chart truth. The implemented experiment crosses context and uncertainty policy, adds corrupted-context and deterministic-copy controls, and preserves declared evaluation denominators after interruptions. No live language-model experiment has been run; H1--H3 remain untested. The checkpoint contributes an auditable diagnostic instrument and a controlled evaluation design, with explicit limits on label independence and interpretation of tool-assisted accuracy.
\end{abstract}
\noindent\textbf{Code:} \url{https://github.com/bothermeQAQ/DSC180A-BaZiReason}\\
\textbf{Audited checkpoint:} November 3, 2026. Calendar diagnostics only.

\section{Introduction}
\subsection{Research problem and scope}
A calendar answer can be fluent and internally plausible while selecting the wrong reference clock or transition rule. BaZi construction makes these choices observable: a chart contains four ordered pairs for the year, month, day, and two-hour period. Each combines one of ten Heavenly Stems and one of twelve Earthly Branches in a sixty-position cycle. In the implementation profiles studied here, the year changes at the Lichun solar-term instant and months change at twelve alternating solar terms, called Jie. These boundaries follow solar motion rather than Gregorian month starts or Lunar New Year \citep{hko2025,lunar}.

Small input changes can cause discrete output changes. A time near 23:00 raises a day-boundary convention question. A repeated daylight-saving wall time can represent two different instants; Python represents that ambiguity using the fold attribute \citep{zoneinfo}. A missing time, zone, or convention may prevent a unique answer. Reliability therefore involves deciding when an answer is supported as well as calculating pillars correctly.

The central question is: \emph{Can deterministic context improve the reliability of convention-specified Chinese calendar answers, and what explains any improvement?} Correctness means agreement with a named profile and, where available, external calendar evidence. It does not establish that a chart predicts personality, health, or future events. Nor does a correct answer copied from a tool establish improved internal reasoning.

Three prospective hypotheses organize the study. H1 predicts more direct-model errors on boundary cases than on ordinary dates. H2 predicts higher exact pillar agreement when the same model receives a complete deterministic result; controls assess whether this amounts to copying or to useful preprocessing. H3 predicts fewer unsupported definitive answers when uncertainty policy is explicit. All remain untested by live models. The checkpoint establishes software and calendar evidence, not language-model performance.

\subsection{Related literature and prior work}
PAL delegates execution to an interpreter after a model constructs an intermediate program \citep{pal}. Toolformer instead trains models to select and incorporate API calls \citep{toolformer}. BaZiReason implements fixed application-side context, without generated programs, learned API selection, or autonomous function calling. Those methods motivate separating language generation from computation, but their published gains do not transfer automatically to this task.

TimeBench organizes temporal reasoning into multiple phenomena and compares model behavior across categories \citep{timebench}. That motivates stratified reporting. It does not establish that the particular boundaries selected here cause errors; H1 requires observations on these inputs and control of task composition.

BaZi-specific prior work substantially overlaps with the proposed setting. BaziQA-Benchmark evaluates contest questions conditioned on fixed chart contexts, separating downstream symbolic judgments from chart derivation \citep{baziqa}. The paper uses 200 questions from the 2021--2025 competitions; its current repository has expanded and revised data. Those versions should not be conflated. Contest answer keys permit scoring but do not establish the scientific validity of astrological prediction. This project uses neither contest questions nor personal birth profiles.

The Deterministic Bazi Calendar Rules Benchmark already targets chart computation, Zi-hour and solar-boundary errors \citep{detbenchmark}. At the inspected commit, its solar solver and model runner raise an unimplemented-operation exception, and its README labels performance estimates provisional. Those estimates are not used as verified evidence. Current BaziQA chart code also exposes time conventions and uses lunar-python, so convention awareness and BaZi benchmarking are not novel in themselves.

The feasible contribution is narrower: executable boundary and missing-information diagnostics, explicit provenance, matched inference conditions, controls for answer transmission and faulty tool context, and separate uncertainty-policy scoring. The study becomes meaningful LLM research only when actual models are evaluated under these contrasts. The present deliverable is its instrument and research design.

\subsection{Data, computation, and evaluation design}
\paragraph{Versioned synthetic data.}
The initial 96-case artifact is preserved in Git at tag \nolinkurl{audit/pre-review-a50f967}. Independent regeneration and a fresh installation reproduced its 42 tests. Inspection nevertheless found four identical Shanghai inputs among its timezone rows, only six of twelve Jie, and same-date 00:00 examples that were not adjacent to the preceding 23:00 examples. Test success therefore did not establish sufficient design coverage.

The revised suite, generated with seed 1802026, contains 152 original synthetic records (Table~\ref{tab:data}). There are 126 scorable computations, 16 ambiguous requests, six unsupported requests and four invalid inputs. Exact input/profile pairs are unique. Forty-eight cases bracket all twelve Jie in 2024 and 2025 by two minutes. Zi-hour groups now include successive 23:59:59 and 00:00:00 observations under both profiles. Timezone pairs represent distinct instants between groups. Ten new matched groups contain 24 omitted-versus-specified cases for convention, time, zone or DST fold. Their 14 scorable counterparts prevent ambiguity evaluation from rewarding indiscriminate refusal.

\begin{table}[ht]
\centering
\caption{Revised suite. Counts describe diagnostic coverage, not independent population samples.}
\label{tab:data}
\begin{tabular}{lr@{\hspace{1em}}lr}
\toprule
Category & Count & Category & Count\\
\midrule
Ordinary noon dates & 20 & Jie boundary pairs & 48\\
New-year contrasts & 8 & Zi/midnight profiles & 20\\
Leap days & 4 & Equivalent time zones & 8\\
DST gaps/folds & 8 & Legacy policy cases & 12\\
Matched H3 cases & 24 & Total & 152\\
\bottomrule
\end{tabular}
\end{table}

Each JSONL record specifies an input, selected profile, category, group ID, expected status and optional pillars, implementation version, validation status, and boundary metadata. Groups retain related transition sides, timezone equivalents, profile variants and missing-field counterparts. Equivalent normalized inputs do not cross groups. Future splits and resampling must preserve those groups. The sample is deliberately boundary enriched and cannot estimate natural request prevalence. Ordinary dates are not matched controls for every transition, so H1 initially supports a descriptive comparison rather than a causal difficulty claim.

\paragraph{Computation contract.}
The wrapper uses lunar-python 1.4.8 \citep{lunar}. It interprets a naive local wall time with IANA zone and optional fold, validates civil-time round trips, and converts the instant to fixed UTC+08:00 for \emph{all} pillars. This reference-clock task does not implement arbitrary birthplace-local charts. The admitted reference-year range is 1901--2099, a study restriction rather than a validated guarantee of the dependency's entire range. True solar time, local civil-clock charts, lunar-month boundaries and Lunar New Year year-boundary profiles are explicitly unsupported.

Sect 1 advances the day at 23:00; sect 2 retains the civil-date day until midnight. Both retain the library's 23:00-advanced hour stem. Deriving sect-2 hours from the displayed day would change this implementation's rule. Missing profiles require clarification, and repeated wall times require fold. Audit fixes load timezone bytes directly from the pinned tzdata package, avoiding the host constructor cache; extreme out-of-range conversions now return a controlled status instead of raising an overflow exception. These are scoped implementation choices, not universal calendrical assertions.

\paragraph{Reference evidence and its limits.}
Expected pillars originate from lunar-python, and every label retains single-engine provenance. Replaying stored labels through that engine is an integrity check, not validation. The historical 184 agreeing sxtwl cells comprised 80 day, 52 year and 52 month checks, reducing to 117 distinct reference-date/pillar queries. All hours and solar-term-date year/month comparisons were excluded. The apparent sample size must not be interpreted as 184 independent validated charts.

The revised suite has 270 agreeing day/year/month cells but only 169 distinct reference queries. The date-level comparator still excludes intra-day year/month crossings. A separately reported, explicitly convention-matched sxtwl hour comparison agrees on 126 cases. Both libraries share ShouXing astronomical lineage \citep{lunar,sxtwl}; such agreement cannot exclude shared errors. Two published HKO day labels and a minute-resolution Lichun value provide small external anchors \citep{hko2025}, without establishing complete-chart correctness.

A distinct numerical path now checks the disputed timestamps. Using Skyfield 1.53, JPL DE421, apparent geocentric solar longitude in the ecliptic of date, and bundled time-scale conversion, the audit computes all 24 Jie in 2024--2025 \citep{skyfield}. Relative to those results, lunar-python differs by $-1.168$ to $+2.261$ seconds and sxtwl by $-15.783$ to $-13.019$ seconds. Versions, kernel hash, definitions and every comparison are retained. This supports lunar-python's timestamps more closely for the selected years; it neither proves absolute accuracy nor diagnoses the libraries' discrepancy. No guessed correction or silent relabeling is applied. The two-minute benchmark offsets lie outside these observed discrepancies, while second-level tests assert only the selected engine's contract.

\paragraph{H2 controls: transmission, preprocessing and faulty context.}
The direct condition receives input, profile, common computation rules and a structured schema, without labels or boundary annotations. The full-context condition additionally receives the complete computed answer. A deterministic copy baseline returns that context without any model call. Its executed 126/126 pillar agreement is tautological: it demonstrates why full-context accuracy alone cannot identify reasoning improvement.

Two additional conditions address this limitation. Timestamp-only context supplies normalization and provenance without pillars, output status or reasons, isolating preprocessing assistance. Corrupted context changes one scorable hour stem to another valid cycle pair while retaining its branch; all other fields and instructions match full context. Exact accuracy and erroneous-hour following measure this specific reliance mechanism. They do not reveal internal reasoning or general robustness to every tool error. The preprocessing output can also reveal that normalization failed, which is part of that condition's information advantage.

\paragraph{H3 policy ablation.}
Each model context condition has explicit and minimal policy versions. Both retain identical input, selected convention, calendar/range rules and output schema. Only the explicit version instructs clarification for missing information and rejection for invalid or unsupported inputs. Thus the factor changes when to abstain, rather than withholding calendrical knowledge. The primary H3 contrast is direct explicit minus direct minimal in definitive-answer frequency on the 16 ambiguous cases; a negative difference is favorable. Full-context policy pairs are secondary because their supplied non-answer is an easier task.

Definitive-answer frequency is reported alongside correct ambiguity handling, valid-output rate and coverage on scorable inputs. Missing responses count as no definitive claim but also as policy/coverage failures; silence is not credited as successful uncertainty handling. Reasons are checked for nonblank structure, not semantic helpfulness. More diverse wording and independent clarification adjudication remain final-project work.

\paragraph{Measurement and reproducibility.}
Exact accuracy uses the 126 labeled observations per full trial. Joint exact accuracy additionally requires correct profile metadata. Per-pillar accuracy uses the same labeled denominator; invalid and missing outputs fail. Non-answer policy metrics use the 26 non-answer cases, with ambiguity reported separately. Repeated consistency compares complete valid responses, without equating consistency with correctness.

The runner holds model, decoding settings, schema, attempts and cases constant, rotates condition order, and records requested/served model, prompt version/hash, timestamp, returned usage, latency and failures. A planned manifest is written before inference. Evaluation enforces its dataset hash, case IDs, methods and trials, retaining interruptions in denominators. Duplicate JSON keys and aliased trial identities are rejected. Descriptive paired intervals resample whole group IDs; small purposive samples and repeated related observations still limit interpretation. Monetary cost remains unknown unless measured.

\paragraph{Research status and feasibility.}
The audit independently reproduced the original software evidence and added regression, coverage and control tests, including a fresh Python 3.12 installation. Calendar results, mocks and deterministic controls are separately labeled. No configured API or usable local model was found in the inspected environment, so no live model experiment or H1--H3 result is reported. A representative 22-case selection supports a bounded pilot; it is a debugging sample, not a powered hypothesis test.

Before the December final, priorities are a real same-model pilot, a frozen grouped evaluation set, expansion of matched ordinary/boundary comparisons, repeated trials and error analysis. Parameter-efficient fine-tuning is deferred. The main remaining limitations are incomplete independent chart evidence, restricted profiles, contamination risks from public calendar sources, purposive and correlated sampling, provider drift, and the inability to infer internal reasoning from output accuracy alone.

\clearpage
\bibliographystyle{plainnat}
\bibliography{references}
\end{document}
